\section{A weighted affine law}\label{sec:law}

We next turn excessive operator growth into an exact scaling identity
for a probability law. For the remainder of the argument suppose,
toward a contradiction, that
\begin{equation}\label{eq:contrary}
 \beta\ge1/4.
\end{equation}
All laws below are constructed at this fixed exponent.

\begin{proposition}[Exact-parameter affine law]\label{prop:law}
Assume \(\beta=\log_2\rho\ge1/4\), where \(\rho\) is defined in
\eqref{eq:rho}. There is a Borel probability \(P\) on \(\KK\) such that, writing
\(\E\) for expectation under \(P\),
\begin{equation}\label{eq:law}
 \E\bigl[|Q_F(x+iy)|\,\Phi(T_{x+iy}F)\bigr]
 =y^{-\beta}\E\Phi(F)
\end{equation}
for every \(x+iy\in\HH\) and every nonnegative Borel test
\(\Phi:\KK\to[0,\infty]\).
\end{proposition}
\begin{proof}
We use the representation of positive functionals on continuous
functions of a compact metric space by Borel measures
\cite[Theorem 7.2]{Folland}.
Compactness of probabilities follows by taking a diagonal limit on
a countable dense family of continuous tests and applying this
representation theorem.

\emph{The operator eigenmeasure.} First construct a probability \(\mu\) with \(\mu L=\rho\mu\),
where measures act on the left as integration functionals.
For \(\tau>\rho\), the series
\[
 R_\tau=\sum_{n\ge0}\tau^{-n}L^n,\qquad
 g_\tau=R_\tau1,\qquad M_\tau=\max g_\tau
\]
converges in operator norm. Since
\[
 Lg_\tau=\tau(g_\tau-1)
 \le\tau(1-1/M_\tau)g_\tau,\qquad 1\le g_\tau\le M_\tau,
\]
positivity gives
\(\|L^n\|\le M_\tau[\tau(1-1/M_\tau)]^n\).
Thus \(\rho\le\tau(1-1/M_\tau)\), so
\(M_\tau\ge\tau/(\tau-\rho)\to\infty\).
At a maximizer \(F_\tau\) of \(g_\tau\), the functional
\(\mu_\tau=\delta_{F_\tau}R_\tau/M_\tau\) is a probability and
\[
 \mu_\tau L=\tau\mu_\tau-\frac{\tau}{M_\tau}\delta_{F_\tau}.
\]
A weak limit as \(\tau\downarrow\rho\) gives the required \(\mu\).

\emph{Horizontal and scale averages.} The composition law gives
\(A_{u+iv}A_{2+i}=A_{u+2v+iv}\): right translation by \(2\)
shifts the grid in \eqref{eq:dyadic} by one cell of width \(2v\).
Using \(\mu=\rho^{-n}\mu L^n\), all but the two endpoint terms cancel:
\[
 \mu A_{2+i}-\mu
 =\rho^{-n}v\bigl(\mu A_{1+v+iv}-\mu A_{-1+v+iv}\bigr).
\]
The norm of this difference is at most \(2C\rho^{-n}\) by
\eqref{eq:strip}, and tends to zero because \(\beta>0\).
Therefore \(\mu A_{2+i}=\mu\).
It follows that \(t\mapsto\mu A_{t+ih}\) has period two for every
fixed \(h>0\).

Define
\(\bar\mu=\tfrac12\int_0^2\mu A_{x+i}\,dx\) on continuous tests.
Periodicity gives \(\bar\mu A_{s+i}=\bar\mu\) for all \(s\in\R\).
Moreover,
\[
 \rho\bar\mu
 =\tfrac12\int_0^2\mu L A_{x+i}\,dx
 =\tfrac12\int_{-1/2}^{3/2}\mu A_{t+i/2}\,dt
 =\bar\mu A_{i/2}.
\]
The middle equality uses \(t=-1/2+x/2\) and \(t=1/2+x/2\)
for the two summands of \(L\); the last uses period two.
Consequently \(h^\beta\bar\mu A_{ih}\) is multiplicatively
two-periodic in \(h>0\).
Each of these functionals is horizontally invariant on the right,
because \(A_{ih}A_{s+i}=A_{hs+i}A_{ih}\).
The functional
\[
 \nu=\int_1^2 h^\beta\bar\mu A_{ih}\,\frac{dh}{h}
\]
has finite, strictly positive mass by compact-family continuity.
Periodicity and the substitution \(u=hy\) give
\(\nu A_{iy}=y^{-\beta}\nu\). Horizontal invariance persists,
so normalization of \(\nu\) gives \eqref{eq:law} for continuous tests.
For each fixed \(x+iy\), the two sides define finite Borel measures
(one is a weighted pushforward). Equality on continuous tests makes
them equal as measures, proving the assertion for all nonnegative
Borel tests.
\end{proof}

The law now supplies the differential moments needed for the signed
Jacobian comparison. For \(z=x+iy\), define
\begin{equation}\label{eq:qk}
 q_F(z)=iy\,\frac{Q_F'(z)}{Q_F(z)}
       =-iy\,\frac{F''(z)}{F'(z)}=a_F(z)+ib_F(z),
 \qquad k=\frac{\beta+2}{3}.
\end{equation}
Thus \(3(k-1)=\beta-1\) and \(3/4\le k\le1\).
Equation~\eqref{eq:action} implies
\(q_{T_zF}(w)=q_F(z\cdot w)\).

\begin{lemma}[Logarithmic-derivative moments]\label{lem:moments}
The law satisfies
\begin{equation}\label{eq:moments}
 \E a_F(i)=-\beta,\qquad
 \E|q_F(i)|^2=\beta(\beta+1).
\end{equation}
\end{lemma}
\begin{proof}
The constant test in \eqref{eq:law} gives
\(\E|Q_F(x+iy)|=y^{-\beta}\).
Since \(\log|Q_F|\) is harmonic,
\[
 \partial_y|Q_F|=\frac{|Q_F|a_F}{y},\qquad
 \Delta|Q_F|=|Q_F|\left|\frac{Q_F'}{Q_F}\right|^2.
\]
Differentiate at \(i\), where \(Q_F(i)=1\).
Lemma~\ref{lem:compact} and local nonvanishing give the uniform
smoothness bounds needed to differentiate under expectation.
\end{proof}
